\begin{theorem}[Multi-agent \(\Sfive\)]
For every premise set \(\Gamma\) and formula \(A\),
\(\Gamma\derives_{\Sfive_I}A\) iff \(A\) is a local consequence of
\(\Gamma\) over epistemic frames in which every \(\sim_i\) is an equivalence
relation.
\end{theorem}

\paragraph{Proof sketch.}
Soundness checks each agent-indexed axiom against
reflexivity, transitivity, and Euclideanness of \(\sim_i\); necessitation is
theorem-only.  For completeness, take maximally consistent sets as worlds and
define \(w\sim_i v\) iff every \(K_iA\in w\) has \(A\in v\).  The
truth lemma is a structural induction.  The three non-distribution schemas
make each canonical relation reflexive, transitive, and Euclidean, hence an
equivalence relation; no step relates \(\sim_i\) to \(\sim_j\) for
\(i\ne j\).  Lindenbaum extension, the modal existence lemma, and the full
truth-lemma calculation are omitted; the one-agent instances are controlled by
\cite[Chapter 4]{BlackburnDeRijkeVenema2001}.

For a nonempty finite group \(G\subseteq I\), the language further adds the
operators \(E_G\) and \(C_G\), with \(E_GA\) and \(C_GA\) formulas
whenever \(A\) is a formula.  Everyone-knows is definable by
\[
  E_GA\defeq\bigwedge_{i\in G}K_iA.
\]
Common knowledge \(C_GA\) receives a satisfaction clause of its own:
\(C_GA\) is true at \(w\) iff \(A\) is true at every world reachable
from \(w\) by a path of length \(n\geq0\) whose edges belong to
\(\bigcup_{i\in G}\sim_i\); the length-zero path makes \(w\) itself
reachable, so the clause is reflexive.  Common knowledge is therefore a
reflexive-transitive closure modality; it is not a finite iteration of
\(E_G\) bounded in advance.
